\scanpage{315}

\mspage{54}
% archive sheet number: 314

Si l'objet de $\operatorname{Nerf}_2$ est
\diagram{\begin{tikzcd}[ampersand replacement=\&,column sep=normal]
x\arrow[r,phantom,"{\overset{f}{\unsure{\text{ --- }|\longrightarrow}}}"{description}] \&
y\arrow[r,dash,"g"] \& z,
\end{tikzcd}}
la construction précédente nous dit, en termes de la donnée combinatoire i.e. de la suite de \deleted{\unsure{données}}
\[
(\varepsilon_1,\ell_1,\pi_1),\ldots,(\varepsilon_n,\ell_n,\pi_n),
\]
quels sont $x,y,z$ : chacun d'eux est le produit dans $\deleted{\unsure{\operatorname{Ob}}}\operatorname{Nerf}_0=\mathrm{Ty}$ des $n$ facteurs $x_i,y_i$\inserted{\unsure{ou|et|,}} $z_i$ respectivement, pour $i\in[1,n]$, donnés dans la définition des $f_{(\varepsilon_i,\ell_i,\pi_i)}$, i.e. de
\[
\Delta_{\varepsilon_i},\quad\Delta_{\varepsilon_i}^{\ell_i},\quad\text{et}\quad
\Delta_{\varepsilon_i}^{\pi_i(1)}\Delta_{\varepsilon_i}^{\pi_i(\unsure{1|2})}\cdots
\Delta_{\varepsilon_i}^{\pi_i(\ell_i)}=\Delta_{\varepsilon_i}^{\operatorname{long}(\pi_i)}
\]
(où $\operatorname{long}\pi_i=\sum_{j\in[1,\ell_i}\pi_i(j)$). Donc
\[
\left\{\begin{aligned}
x&=\prod_{i\in[1,n]}^{\mathrm{Ty}}\Delta_{\varepsilon_i}\\
y&=\prod_{i\in[1,n]}^{\mathrm{Ty}}\Delta_{\varepsilon_i}^{\ell_i}
\deleted{\unsure{=\Delta_{\varepsilon_i}^{\operatorname{long}(\ell)}}}\\
\unsure{Z|z}&=\prod_{i\in[1,n]}^{\mathrm{Ty}}\Delta_{\varepsilon_i}^{\operatorname{long}(\pi_i)}
\deleted{\unsure{=\Delta_{\varepsilon_i}^{\sum_{i\in[1,n]}\operatorname{long}(\pi_i)}}},
\end{aligned}\right.
\]
\margin{attention, les produits doivent être pris en respectant l'ordre des facteurs donné par la \unsure{notation} indicielle en $i$}
et on trouve de même $f,g$ comme des produits (au sens $*$) de $n$ facteurs $f_i$, resp. $g_i$ $(i\in[1,n])$ explicités dans la formule (1)(2) (p. 52, 53) :
\[
\begin{aligned}
\unsure{\{|[} \qquad f_i&=f_{\varepsilon_i,\ell_i}:\Delta_{\varepsilon_i}\longrightarrow\Delta_{\varepsilon_i}^{\ell_i}\\
g_i&=g_{\varepsilon_i,\pi_i}:\Delta_{\varepsilon_i}^{\ell_i}\longrightarrow\Delta_{\varepsilon_i}^{\operatorname{long}(\pi_i)}\ ; 
\end{aligned}
\]
$g_i$ lui-même étant représenté comme un $*$-produit de $\ell_i$ facteurs \inserted{indécomposables}, de sorte que \inserted{$f$ et} $g$ apparaissent sous forme
